\chapter*{Engineering}
\phantomsection
\addcontentsline{toc}{chapter}{Engineering}

Dieser Teil legt dar, wie Optimierungsergebnisse unter ausdrücklichen
Engineering-Regeln zu zulässigen Eisenbahnlösungen werden.

\section*{Was wissen wir bereits?}
Die vorangehenden Teile haben die mathematischen, geometrischen, operativen,
realisierungsbewussten und optimierungsbezogenen Grundlagen von AIM geschaffen.
Alignments können durch strukturierte operatorbasierte Verfahren dargestellt,
bewertet, analysiert und verbessert werden. Optimierung erzeugt nach
festgelegten Zielen verbesserte Lösungen. Sie allein bestimmt jedoch nicht, ob
eine Lösung ingenieurtechnisch akzeptabel ist.

\section*{Was fehlt noch?}
Eisenbahningenieurwesen arbeitet in einem Rahmen aus Regeln, Empfehlungen,
Normen, Vorschriften und operativen Bedingungen. Nicht jedes mathematisch
machbare Alignment ist zulässig. Gleisgeometrie muss Engineering-Grenzen
einhalten, operative Zustände sicher, Baulösungen praktisch und Alternativen
mit anwendbaren Normen vereinbar bleiben.

Es fehlt daher die Engineering-Bedingungsebene, die mathematische Modelle mit
der Eisenbahnpraxis verbindet.

\section*{Was folgt?}
Dieser Teil führt Engineering-Bedingungen als ausdrückliche AIM-Bestandteile
ein. Normen werden nicht erst nach dem Entwurf konsultiert, sondern als
strukturierte Bedingungen interpretiert, die unmittelbar an Analyse und
Optimierung teilnehmen können.

Die Kapitel untersuchen Normen und Vorschriften, Zulässigkeitsbedingungen,
eisenbahnspezifische Engineering-Regeln und formale Darstellungen operativer
Anforderungen. Ziel ist, Eisenbahnnormen von passivem Referenzmaterial in
aktive Bestandteile der Informationsmodellierung zu überführen.

\section*{Bezug zur AIM-Roadmap}
In der AIM-Roadmap entspricht dieser Teil dem Übergang von Optimierung zu
ingenieurtechnischer Zulässigkeit. Die Skizze beantwortet: Welche
Bedingungsebenen überführen mathematisch verbesserte in ingenieurtechnisch
gültige Zustände?

\begin{center}
\begin{tikzpicture}[node distance=1.4cm,box/.style={draw,rounded corners,minimum width=5cm,minimum height=.9cm,align=center}]
\node[box] (optimization) {Optimierung};
\node[box,below=of optimization] (constraints) {Engineering-Bedingungen};
\node[box,below=of constraints] (rules) {Eisenbahnregeln};
\draw[->,thick] (optimization)--(constraints); \draw[->,thick] (constraints)--(rules);
\end{tikzpicture}
\end{center}

Die Abfolge
\[
\text{Optimierung}\rightarrow\text{Engineering-Bedingungen}
\rightarrow\text{Eisenbahnregeln}
\]
markiert den Übergang von verbesserten Lösungen zur Sicherung ihrer
Zulässigkeit in realen Eisenbahnsystemen. Engineering verbindet mathematische
Optimierung mit praktischer Umsetzung.

\begin{principle}[Prinzip der Engineering-Bedingungen]
\label{princ:engineering-constraint-principle}
In AIM sollen Eisenbahnnormen, Vorschriften und operative Regeln als
ausdrückliche Engineering-Bedingungen dargestellt werden, die unmittelbar an
Modellierung, Analyse und Optimierung teilnehmen können.
\end{principle}

\begin{summary}
Optimierung erklärt, wie Eisenbahnsysteme verbessert werden können.
Engineering erklärt, welche Verbesserungen zulässig sind. Dieser Teil
begründet die Bedingungsebene zwischen mathematischer Optimierung und
praktischem Eisenbahnentwurf und -betrieb.
\end{summary}

Diese Ebene bereitet den Teil Anwendungen vor, in dem zulässige Begriffe in
Interoperabilitäts- und Arbeitsablaufkontexten erprobt werden.
